% =====================================================================
%  Chapter 3: Data, Measurement, and Open Infrastructure
%  Source: P3 transcription (ACS, GDP/GDI, price indices) + new content
% =====================================================================

\chapter{Data, Measurement, and Open Infrastructure}
\label{chap_Data}

\section{The Data Landscape for Urban and Regional Research}
\label{sec_DataLandscape}

Modern urban and regional research rests on an unusually rich base of public
data infrastructure. The US federal statistical system produces consistent,
publicly accessible data on population, housing, labour markets, income,
prices, and the built environment at fine geographic resolution---often down
to the census block or census tract. This infrastructure is not merely a
convenience; it is the empirical foundation that makes regional policy
informatics possible without proprietary datasets or expensive licensing fees.

This chapter provides a practical guide to the key data sources, their
structure and limitations, and the open-source R and Python tools for
accessing them. It also covers measurement concepts---price indices, national
accounts, margins of error---that recur throughout the book.

\noindent\textbf{A note on geographic aggregation.} Before working with any
spatial data, it is essential to understand the areal unit at which
observations are recorded and the implications of aggregation choices. The
Modifiable Areal Unit Problem (MAUP)\index{MAUP|see{Modifiable Areal Unit
Problem}}\index{Modifiable Areal Unit Problem} arises because statistical
results can change substantially when the same data are aggregated to
different zone systems---from census tracts to ZIP codes to counties to
metropolitan areas. The MAUP is treated at length in
Chapter~\ref{chap_Spatial}; here we note only that the choice of geographic
unit is never purely technical, and results should be routinely checked for
sensitivity to that choice.


\section{Survey Data and Statistical Precision}
\label{sec_SurveyData}

\subsection{American Community Survey (ACS) and Margins of Error}

ACS tables report \textbf{margins of error} (MOE) that permit construction
of 90\% confidence bands around estimates. The relationship between MOE
and standard error is:

\begin{tcolorbox}[keyresult]
\begin{equation}
  \text{s.e.} = \frac{\text{MOE}_\text{ACS}}{1.645}
  \qquad\Longleftrightarrow\qquad
  \text{MOE} = \text{s.e.} \times z\text{-score}
\end{equation}
\end{tcolorbox}

\noindent\textbf{Testing differences in population means.} When comparing
an estimate $M_1$ from geography 1 with estimate $M_2$ from geography 2,
using standard errors $s_1$ and $s_2$ derived from the ACS MOEs:
\begin{equation}
  \left|\frac{M_1 - M_2}{\sqrt{s_1^2 + s_2^2}}\right| > z_\text{confidence}
\end{equation}

If this test statistic exceeds the critical value (e.g.\ 1.645 for 90\%
confidence), the difference is statistically significant at that level.

The ACS is available at multiple geographic levels and in 1-year and 5-year
estimates. The 5-year estimates are more precise (lower MOE) but cover a
longer time window; the 1-year estimates are more timely but available only
for geographies with population $\geq 65{,}000$.


\section{National and Regional Economic Accounts}
\label{sec_NationalAccounts}

\subsection{GDP and Gross Domestic Income (GDI)}

Gross Domestic Income (GDI) measures output from the income side
(NIPA)\index{NIPA|see{National Income and Product Accounts}}. Its
components sum to GDP:

\begin{table}[htp]
\centering
\caption{GDI components as shares of GDP (approximate, 2003)}
\label{tbl_GDI}
\renewcommand{\arraystretch}{1.3}
\begin{tabular}{@{}lc@{}}
\toprule
Component & Approximate share \\
\midrule
Compensation of employees & 56\% \\
Taxes on production and imports (TOPI), less subsidies & 7\% \\
Net operating surplus (NOS) & 25\% \\
Consumption of fixed capital (CFC) & 12\% \\
Statistical discrepancy & $\approx 0.2\%$ \\
\midrule
\textbf{Total $\equiv$ GDP} & \textbf{100\%} \\
\bottomrule
\end{tabular}
\renewcommand{\arraystretch}{1}
\end{table}

\noindent where TOPI $\equiv$ sales taxes, property taxes, custom duties;
NOS $\equiv$ corporate profits, net earnings of unincorporated businesses,
government surplus; CFC $\equiv$ depreciation, capital expenditures, interest.

\subsection{Measuring Price Changes}

\begin{itemize}
  \item \textbf{CPI:}\index{CPI|see{Consumer Price Index}}
    Fixed basket of goods tracked over time; historically a
    Laspeyres index.
  \item \textbf{GDP deflator:}
    $\displaystyle\frac{\text{Nominal GDP}}{\text{Real GDP}}\times100$.
    Differs from CPI in that basket weights are \emph{not} fixed.
\end{itemize}

Two classical fixed-weight price indices:

\begin{table}[htp]
\centering
\caption{Fixed-weight price indices and their biases}
\label{tbl_PriceIndices}
\renewcommand{\arraystretch}{1.6}
\begin{tabular}{@{}llc@{}}
\toprule
Index & Formula & Bias \\
\midrule
\textbf{Paasche}\index{Paasche index}
  & $\dfrac{\displaystyle\sum p_{it}q_{it}}{\displaystyle\sum p_{i0}q_{it}}$
  & Understates inflation \\[4pt]
\textbf{Laspeyres}\index{Laspeyres index}
  & $\dfrac{\displaystyle\sum p_{it}q_{i0}}{\displaystyle\sum p_{i0}q_{i0}}$
  & Overstates inflation \\
\bottomrule
\end{tabular}
\renewcommand{\arraystretch}{1}
\end{table}

The BLS now uses the \textbf{Fisher ideal index}\index{Fisher ideal index},
the geometric mean of the Paasche ($P$) and Laspeyres ($L$) indices:
\begin{equation}
  \text{Fisher} = \sqrt{P \times L}
  \qquad\text{cf.}\quad
  \underbrace{\Bigl(\prod_{i=1}^{n}a_i\Bigr)^{1/n}}_{\text{geometric mean}}
  \;\text{vs.}\;
  \underbrace{\frac{1}{n}\sum_{i=1}^{n}a_i}_{\text{arithmetic mean}}
\end{equation}


\section{The Open Data Stack}
\label{sec_OpenData}

Table~\ref{tbl_OpenData} summarises the key public data sources for regional
analytics and the primary R packages for accessing each. Together these
sources enable a complete regional analytics workflow without any proprietary
data.

\begin{table}[htp]
\centering
\caption{Public data infrastructure for regional analytics}
\label{tbl_OpenData}
\renewcommand{\arraystretch}{1.4}
\begin{tabular}{@{}p{3.8cm}p{5.5cm}p{3cm}@{}}
\toprule
\textbf{Source} & \textbf{Content} & \textbf{R package} \\
\midrule
Census / ACS & Demographics, housing, income by geography & \texttt{tidycensus} \\
TIGER/Line & All US administrative geographies & \texttt{tigris} \\
FRED & Macro and regional time series & \texttt{fredr} \\
LEHD / LODES & Workplace--residence flows at block level & \texttt{lehdr} \\
BLS QCEW & Employment by county $\times$ industry $\times$ quarter & Direct API \\
BEA Regional Accounts & GDP by metro; personal income by county & \texttt{bea.R} \\
HUD portal & CHAS, HMDA, LIHTC, Fair Market Rents & \texttt{hudapi} \\
IPUMS & Harmonised microdata 1850--present & \texttt{ipumsr} \\
OpenStreetMap & Networks, amenities, land use & \texttt{osmdata} \\
Zillow Research & ZHVI, ZORI at ZIP code level & Direct download \\
EPA EJSCREEN & Environmental justice by census geography & Direct API \\
USDA ERS & Rural--urban classification; agricultural land & Direct download \\
DOT HPMS & Highway infrastructure & Direct download \\
FTA NTD & Transit system data & Direct download \\
\bottomrule
\end{tabular}
\renewcommand{\arraystretch}{1}
\end{table}

\noindent\textbf{The R stack that replaces ArcGIS.} The combination
\texttt{sf} + \texttt{terra} + \texttt{tidycensus} + \texttt{osmdata} +
\texttt{spdep} + \texttt{spatialreg} covers essentially all ArcGIS
social-science applications, with the crucial difference that every operation
is scriptable, reproducible, and version-controlled. This is a scientific
claim about the validity of methods---not merely a software preference. A
point-and-click GIS workflow produces irreproducible analysis; an R workflow
produces reproducible analysis.


\section{Reproducible Workflows: The Technical Infrastructure}
\label{sec_Reproducible}

The open-source commitment described in the Preface requires a minimum
technical infrastructure. Three tools form the core:

\begin{enumerate}
  \item \textbf{Quarto} (\texttt{quarto.org}): A literate programming
    platform that combines prose, code (R or Python), and mathematical
    notation in a single source file, rendering to PDF, HTML, or
    presentations. Every example in this book was written and rendered
    using Quarto.

  \item \textbf{GitHub}: Version control and collaboration. Maintaining a
    public repository of code and data supporting published analyses is
    increasingly treated as a condition of publication in economics journals
    and should be the norm in planning and regional science.

  \item \textbf{Package management}: \texttt{renv} (R) and
    \texttt{requirements.txt} / \texttt{pyproject.toml} (Python) pin
    package versions to ensure that an analysis runs identically on every
    machine, now and in future. Run \texttt{renv::snapshot()} after
    installing any new R package.
\end{enumerate}

\noindent\textbf{Agentic coding assistance.} Large language model assistants
(GitHub Copilot, Claude, ChatGPT) can generate competent spatial regression
code from a plain-English description. This capability makes the theoretical
content of this book \emph{more} important, not less. When the code is
generated automatically, the limiting factor on analysis quality becomes
whether the analyst can evaluate the output---does the weight matrix choice
make sense? Is the identification strategy credible? Does the coefficient
have a causal interpretation? These are questions that require the material
in Parts~II and~III; they cannot be answered by a coding assistant.

% TODO: Add worked example of tidycensus data pull and ACS MOE computation

\renewcommand\bibname{Further reading}
\begin{subappendices}
\bibliographystyle{econometrica}
\bibliography{Literature/OverLord}
\IfFileExists{Literature/Data.bbl}{\chapter{Data sources}\label{chap_Data}}{\typeout{Need bibtex on Data}}
\end{subappendices}
